% SPDX-License-Identifier: CC-BY-SA-4.0
% Author: Matthieu Perrin
% Part: <Nom de la partie>
% Section: <Nom de la section>
% Sub-section: <Nom de la sous-section>  % (facultatif, laisser vide si non utilisé)
% Frame: <Titre de la slide>

\begingroup

\begin{frame}{Difficulté du problème universel}

  \vspace{-3mm}
  \begin{block}{Définitions -- \textsc{re}-difficulté et \textsc{re}-complétude}
    \vspace{-1mm}
    Soit $A \in \textsc{lang}$ un problème de décision. On dit que : 
    \begin{itemize}
    \item $A$ est \structure{\textsc{re}-difficile} si tout problème de \textsc{re} se réduit à $A$ : 
      $\alert{\forall P \in \textsc{re},~ P \leq_m A}$
      \begin{itemize}
      \item\structure{Intuition.} $A$ est au moins aussi difficile que tous les problèmes de \textsc{re}
      \end{itemize}
    \item $A$ est \structure{\textsc{re}-complet} si $A \in \textsc{re}$ et $A$ est \structure{\textsc{re}-difficile}
      \begin{itemize}
      \item\structure{Intuition.} $A$ est le plus difficile des problèmes de \textsc{re}
      \end{itemize}
    \end{itemize}
    Ainsi, le degré des problèmes \textsc{re}-complets est le maximum de $\textsc{re}$ pour $\le_m$. 
  \end{block}

  \pause
  \begin{block}{Théorème -- \textsc{re}-complétude de \textsc{Universal}}
    \vspace{-1mm}
    Pour tout encodage effectif $\langle \mathcal{L}, \llbracket \cdot \rrbracket \rangle$, \alert{$\textsc{Universal}_{\langle \mathcal{L}, \llbracket \cdot \rrbracket \rangle}$ est \textsc{re}-complet}. 

    \vspace{2mm}
    \structure{Démonstration.} Soit $P \in \textsc{re}$. Il existe une machine $M_P$ qui semi-décide $P$.

    On a la réduction $f$ de $P$ vers $\textsc{Universal}$ suivante :
    $$\alert{f : \left\{\begin{array}{rcl}
      \mathcal{I}(P) & \rightarrow & \mathcal{I}(\textsc{Universal}) \\
      x              & \mapsto     & \langle M_P, x \rangle \\
      \end{array}\right.}$$ 
    
    \begin{itemize}
    \item\vspace{-2mm} $f$ est calculable (transformation syntaxique)
    \item $\forall x \in \mathcal{I}(P), \quad
      \alert{f(x) \in \textsc{pos}(\textsc{Universal})}
      \quad\Leftrightarrow \quad
      x \in \llbracket M_P \rrbracket
      \quad\Leftrightarrow\quad
      \alert{x \in \textsc{pos}(P)}$
    \end{itemize}
  \end{block}

\end{frame}

\endgroup
